% \iffalse meta-comment
%
% Copyright (C) 2017- by Yanshuo Chu <yanshuoc@gmail.com>
%
% This file may be distributed and/or modified under the
% conditions of the LaTeX Project Public License, either version 1.3a
% of this license or (at your option) any later version.
% The latest version of this license is in:
%
% http://www.latex-project.org/lppl.txt
%
% and version 1.3a or later is part of all distributions of LaTeX
% version 2004/10/01 or later.
%
% \fi
%
% \section{编号后间距与固定名字}
%
% \changes{v3.2a}{2026/09/13}{从 \file{hit-defaults.dtx} 拆出来编进类文件（进不了
%   \file{hithesis.cfg}）；\cs{PGR} 与 \cs{UGR} 两个类共用，书名号挪到链接外}
%
% 两组取值，都进不了 \file{hithesis.cfg}：上面那组是循环生成的，而
% \file{cfg} 里一行一个键；下面那组是 \cs{def}，而 \file{cfg} 只收
% \cs{hit@presetup}。所以它们编进类文件，守卫叫 \texttt{shared-names}。
%
%    \begin{macrocode}
%<*shared-names>
%    \end{macrocode}
%
% 编号之后排什么，默认一律空，空的含义是「用各处原有的取值」，见
% \cs{hit@after@number:nnn}。这十二条不能省：常量是用到才 \cs{cs\_gset} 出来的，
% 不设的话名字根本不存在，取值时会被当成 \cs{relax}、判成非空，间距静悄悄地消失。
%    \begin{macrocode}
\clist_map_inline:nn
  {
    after-chapter-number , after-section-number ,
    after-subsection-number , after-subsubsection-number ,
    after-figure-caption-number , after-table-caption-number
  }
  {
    \hit@presetup{ term = { #1-zh = { } , #1-en = { } } }
  }
%    \end{macrocode}
%
% \cs{hithesis}、\cs{hit}、\cs{PGR}、\cs{UGR} 是模板与学校的固定名字，不是可配置
% 的词汇，所以直接定义，不进 \texttt{hit/term} 族：进了族用户就能改，改出来的
% 「哈尔滨工业大学」不叫哈尔滨工业大学。
%
% 这一段要退出 \pkg{expl3} 的 catcode 写。\cs{def}\cs{PGR} 的值里有
% \verb|\hit 研究生|，那个空格是控制词的终止符；\pkg{expl3} 下空格被忽略，
% 后面的汉字又是 letter，整串会粘成一个控制序列。
%    \begin{macrocode}
\ExplSyntaxOff
\def\hithesis{\textsc{hi}\-\textsc{Thesis}}
\def\hit{哈尔滨工业大学}
% 书名号搁在链接外头：整个书名装进 \cs{href} 的话，闭书名号“》”会顶着
% 组尾，右半宽的剪切在老版 xeCJK 里过不了组界，跨 TeX Live 版本要翻断点，
% 2026 倒是能跨组压上。号在外面，两边就都按素的“闭号接汉字”排了。
\def\PGR{《\href{http://hitgs.hit.edu.cn/aa/fd/c3425a109309/page.htm}{\hit 研究生学位论文撰写规范}》}
\def\UGR{《\href{http://jwc.hit.edu.cn/2566/list.htm}{\hit 本科生毕业论文撰写规范}》}
\ExplSyntaxOn
%</shared-names>
%    \end{macrocode}
%
% \Finale
